\documentclass[tikz,border=12pt]{standalone}
\usepackage[T1]{fontenc}
\usepackage[utf8]{inputenc}
\usepackage{lmodern}
\usepackage[skins]{tcolorbox}
\usetikzlibrary{arrows.meta}

\definecolor{ink}{HTML}{25344A}
\definecolor{quebg}{HTML}{EEF2F8}
\definecolor{queborder}{HTML}{405B7B}
\definecolor{comobg}{HTML}{EAF6EE}
\definecolor{comoborder}{HTML}{34835A}

% ---- Estilo "minicard": mini-header colorido (sigla, texto branco) +
% corpo branco embaixo (descrição, texto escuro). ----
\tcbset{
  minicard/.style={
    % largura fixa (alinha a grade); altura NÃO é travada (sem "height=") -
    % cresce conforme o conteúdo precisar, como "minimum height" fazia
    % no node original. Travar a altura (height=) foi o que estourava a
    % largura quando o texto não cabia.
    enhanced, arc=4pt, boxrule=0.8pt, colback=white,
    width=2.38cm,
    valign=center, halign=center,
    fonttitle=\bfseries\sffamily\small, coltitle=white,
    halign title=center,
    fontupper=\sffamily\footnotesize, colupper=ink,
    toptitle=3pt, bottomtitle=3pt,
    top=3pt, bottom=3pt, left=3pt, right=3pt,
    titlerule=0pt
  },
  what card/.style={minicard, colframe=queborder!75, colbacktitle=queborder},
  how card/.style={minicard, colframe=comoborder!80, colbacktitle=comoborder}
}

\tikzset{
  relation/.style={
    -{Latex[length=2.5mm]},
    draw=ink!65, line width=0.9pt
  }
}

\begin{document}
\begin{tikzpicture}[font=\sffamily]

  % ---- Faixas de fundo: FIXAS (posição definida nas rodadas anteriores,
  % não mexer mais). ----
  \filldraw[fill=quebg, draw=queborder!70,
            rounded corners=6pt, line width=0.8pt]
    (-1.48,2.865) rectangle (16.73,5.655);

  \filldraw[fill=comobg, draw=comoborder!75,
            rounded corners=6pt, line width=0.8pt]
    (-1.48,-0.465) rectangle (16.73,2.415);

  \node[
    anchor=west, fill=quebg, draw=queborder!70,
    rounded corners=3pt,
    inner xsep=8pt, inner ysep=4pt,
    font=\ttfamily\footnotesize, text=queborder
  ] at (-1.18,5.655)
    {\textbf{CAMADA 1}\quad O QUÊ (CPT)};

  \node[
    anchor=west, fill=comobg, draw=comoborder!75,
    rounded corners=3pt,
    inner xsep=8pt, inner ysep=4pt,
    font=\ttfamily\footnotesize, text=comoborder
  ] at (-1.18,-0.465)
    {\textbf{CAMADA 2}\quad COMO (AIACPT)};

  % ---- Cartões "O QUÊ": mini-header com a sigla + corpo com o assunto ----
  % Corpo em \parbox explícito (2.0cm): não confia mais na quebra de
  % linha "automática" do tcbox, que se mostrou inconsistente.
  \node[inner sep=0pt] (q1) at (0,4.260)
    {\tcbox[what card, title=IC]{\parbox{2.0cm}{\centering Intelligence}}};
  \node[inner sep=0pt] (q2) at (3.05,4.260)
    {\tcbox[what card, title=DS]{\parbox{2.0cm}{\centering Data Services}}};
  \node[inner sep=0pt] (q3) at (6.10,4.260)
    {\tcbox[what card, title=TW]{\parbox{2.0cm}{\centering Trustworthiness (LGPD)}}};
  \node[inner sep=0pt] (q4) at (9.15,4.260)
    {\tcbox[what card, title=IR]{\parbox{2.0cm}{\centering Integration}}};
  \node[inner sep=0pt] (q5) at (12.20,4.260)
    {\tcbox[what card, title=UX]{\parbox{2.0cm}{\centering User Experience}}};
  \node[inner sep=0pt] (q6) at (15.25,4.260)
    {\tcbox[what card, title=MG]{\parbox{2.0cm}{\centering Management}}};

  % ---- Cartões "COMO": mesma lógica, cor verde ----
  \node[inner sep=0pt] (c1) at (0,0.975)
    {\tcbox[how card, title=LA]{\parbox{2.0cm}{\centering Learning and Adaptation}}};
  \node[inner sep=0pt] (c2) at (3.05,0.975)
    {\tcbox[how card, title=PK]{\parbox{2.0cm}{\centering Perception and Knowledge}}};
  \node[inner sep=0pt] (c3) at (6.10,0.975)
    {\tcbox[how card, title=GS]{\parbox{2.0cm}{\centering Governance and Safety (Privacidade)}}};
  \node[inner sep=0pt] (c4) at (9.15,0.975)
    {\tcbox[how card, title=CG]{\parbox{2.0cm}{\centering Cognition and Reasoning}}};
  \node[inner sep=0pt] (c5) at (12.20,0.975)
    {\tcbox[how card, title=AE]{\parbox{2.0cm}{\centering Action and Execution}}};
  \node[inner sep=0pt] (c6) at (15.25,0.975)
    {\tcbox[how card, title=IC]{\parbox{2.0cm}{\centering Interaction and Collaboration}}};

  \foreach \i in {1,...,6}
    \draw[relation] (q\i.south) -- (c\i.north);

\end{tikzpicture}
\end{document}